\documentclass[runningheads]{llncs}

\usepackage[utf8]{inputenc}
\usepackage[T1]{fontenc}
\usepackage{amsmath,amssymb}
\usepackage{booktabs,array,longtable,enumitem}
\usepackage{xcolor}
\usepackage{hyperref}
\hypersetup{
  hidelinks,
  pdftitle={HypG v2.0.1: A Complete Safe-Rust Post-Quantum Cryptographic Suite},
  pdfauthor={HypG Project},
  pdfsubject={Extended technical report on design, formal specification, and evaluation},
  pdfkeywords={post-quantum cryptography, Module-LWE, BLAKE3, ChaCha20, FORS, WOTS+, XMSS, safe Rust}
}
\setlength{\emergencystretch}{5em}
\allowdisplaybreaks


\newcommand{\concat}{\,\Vert\,}
\newcommand{\bits}{\{0,1\}}
\newcommand{\Z}{\mathbb{Z}}
\newcommand{\Ring}{\mathcal{R}_q}
\newcommand{\Hash}{\mathsf{H}}
\newcommand{\XOF}{\mathsf{XOF}}
\newcommand{\KDF}{\mathsf{KDF}}
\newcommand{\PRF}{\mathsf{PRF}}
\newcommand{\ADRS}{\mathrm{ADRS}}
\newcommand{\KeyGen}{\mathsf{KeyGen}}
\newcommand{\Encaps}{\mathsf{Encaps}}
\newcommand{\Decaps}{\mathsf{Decaps}}
\newcommand{\Encode}{\mathsf{Encode}}
\newcommand{\Decode}{\mathsf{Decode}}
\newcommand{\Pack}{\mathsf{Pack}}
\newcommand{\Unpack}{\mathsf{Unpack}}
\newcommand{\CBD}{\mathrm{CBD}}
\newcommand{\Compress}{\mathrm{Compress}}
\newcommand{\Decompress}{\mathrm{Decompress}}
\newcommand{\negl}{\mathsf{negl}}
\newcommand{\ppt}{\mathsf{PPT}}
\newcommand{\adv}{\mathcal{A}}

\newenvironment{steps}[1]
  {\par\medskip\noindent\textbf{#1}\par\nopagebreak\small
   \begin{enumerate}[label=\arabic*:,leftmargin=*]}
  {\end{enumerate}\par\medskip}

\begin{document}
\setcounter{tocdepth}{1}

\title{\textbf{HypG v2.0.1: A Complete Safe-Rust\\
Post-Quantum Cryptographic Suite}\\[0.5ex]
\large Extended Technical Report on\\
Design, Formal Specification, and Evaluation}
\titlerunning{HypG v2.0.1: Extended Technical Report}
\author{HypG Project}
\authorrunning{HypG Project}
\institute{Independent research prototype}
\let\hypgorigaddcontentsline\addcontentsline
\renewcommand{\addcontentsline}[3]{}
\maketitle
\let\addcontentsline\hypgorigaddcontentsline

\begin{abstract}
We present HypG v2.0.1, a from-scratch, zero-external-dependency
post-quantum cryptographic suite written in Rust with
\texttt{\#![forbid(unsafe\_code)]} in every library.  The suite contains
a Module-LWE key-encapsulation mechanism at module rank four over
$\Z_{4096}[X]/(X^{256}+1)$; a portable BLAKE3 hash, XOF, keyed-hash, and
derive-key layer; a ChaCha20 plus keyed-BLAKE3 Encrypt-then-MAC
authenticated-encryption scheme with byte-identical one-shot and
streaming interfaces; and a stateless hash-based signature scheme with
FORS $35\times 9$, WOTS+ $w=16$, and a 17-layer XMSS hypertree.  The
signature wire format carries the separate \textbf{2.0.1} revision:
explicit one-byte $\PRF$/$F$ domain tags and FIPS-205-style promotion of
odd WOTS+ L-tree nodes.

The paper is deliberately both a design document and a normative
technical report.  It gives the exact algorithms and wire layouts,
formal definitions and correctness lemmas, the complete domain registry,
the public API contracts, the frozen known-answer-test inventory, and a
controlled evaluation.  The recorded release has 146 deterministic
tests passing in debug, release, and on the Rust 1.70 MSRV; four frozen
KAT files totalling 234,114 bytes; a 14-suite adversarial harness; a
20,000-trial failure sweep; a full 11,520-bit sweep of one test
ciphertext; Miri coverage of the KEM, AEAD, and DS bounds tests; and
byte-identical KAT output on x86-64 and \texttt{qemu-aarch64}.  The
White Paper companion gives the design rationale and threats-to-validity
discussion; the Yellow Paper gives the same normative material in a
standalone specification format.

None of this is a security proof or an audit.  The KEM's exact FO-style
transform is not covered by a published theorem; the lattice hardness of
the chosen parameters has not been confirmed by an external estimator;
the signature composition has no exact signing-volume bound; the AEAD
composition has no reduction for its exact instantiation; and there is
no external audit, formal verification, \texttt{dudect}-grade timing
study, or power/EM/fault campaign.  HypG is a research prototype and
must not be deployed to protect real assets.
\end{abstract}

\keywords{post-quantum cryptography \and Module-LWE \and
Fujisaki--Okamoto \and hash-based signatures \and FORS \and WOTS+ \and
XMSS \and BLAKE3 \and ChaCha20 \and authenticated encryption \and
safe Rust \and reproducible evaluation.}

\tableofcontents
\newpage

%======================================================================
\section{Introduction}
\label{sec:intro}

Public-key cryptography deployed today is vulnerable to a
cryptographically relevant quantum computer: Shor's algorithm solves
integer factorisation and discrete logarithms in polynomial time.
Symmetric primitives are less affected, but their effective security
halves under Grover search.  A post-quantum suite must therefore combine
quantum-resistant public-key constructions with conservative symmetric
cryptography and a disciplined implementation.

HypG is an experiment in building the \emph{whole} suite rather than a
single primitive: key establishment, hashing, authenticated encryption,
and signatures, sharing one small implementation core.  The project
imposes unusually strict engineering constraints: no external
dependencies, no \texttt{unsafe} code, \texttt{no\_std} cryptographic
libraries, fixed public loop bounds, explicit little-endian encodings,
deterministic test vectors, and mechanically checked CI gates.

The current release is \textbf{v2.0.0} with a \textbf{DS v2.0.1} wire
revision.  It is a research prototype.  The report deliberately
separates (i) what the code does, (ii) what the test suite establishes,
and (iii) what would still be required for a security claim.  This
extended report is a companion to the White and Yellow Papers; it is not
a peer-reviewed publication, an audit, or a certification.

\subsection{Contributions}

The contributions are:

\begin{enumerate}[leftmargin=*]
\item \textbf{A complete, dependency-free implementation.}  A
  \texttt{no\_std}, no-allocation cryptographic core and four protocol
  pillars, with all dependency edges confined to workspace path
  dependencies and all library crates under
  \texttt{\#![forbid(unsafe\_code)]}.
\item \textbf{A precise v2.0 parameter profile.}  A rank-four
  Module-LWE KEM at lattice dimension 1024 with a power-of-two modulus
  $q=4096$, and an SLH-DSA-256f-shaped signature profile with
  $n=32$, $h=68$, $d=17$, $h'=4$, FORS $(35,9)$, and WOTS+ $w=16$.
\item \textbf{A documented DS 2.0.1 revision.}  Explicit one-byte
  $\PRF$/$F$ tags remove a structural collision between the two
  64-byte keyed-hash inputs, and odd WOTS+ L-tree nodes are promoted
  unchanged in the FIPS 205 style.
\item \textbf{An honest treatment of the CA-PRF.}  The cellular
  automaton is presented as a branchless mixing and domain-separation
  stage whose uniformity argument rests on the final BLAKE3 XOF, not
  as an independently proven hardness assumption.
\item \textbf{A corrected KEM correctness analysis.}  The compression
  residuals from $c_1$ and $c_2$ are included explicitly, yielding a
  working standard deviation of approximately 70 and a margin of
  approximately $14.6\sigma$ against the 1024 decision boundary.  This
  is a heuristic estimate, not a formal failure bound.
\item \textbf{A byte-identical streaming AEAD.}  The streaming front
  uses the same subkeys, MAC input, counter schedule, and tag as the
  one-shot construction; only the scheduling differs.
\item \textbf{A reproducible evaluation record.}  Frozen KAT checksums,
  cross-architecture equality, a deterministic red-team harness, a
  local lattice lab, Miri subsets, and a controlled release benchmark
  are recorded with commands and machine details.
\item \textbf{A complete normative specification.}  The report includes
  the exact algorithms, byte layouts, domain registry, API contracts,
  and KAT inventory, not merely a high-level design narrative.
\item \textbf{Explicit non-claims and an audit agenda.}  The report names
  the exact proof, estimator, side-channel, and operational gaps that
  stand between the current artifact and a deployable system.
\end{enumerate}

\subsection{What this report does not claim}

This report does \emph{not} claim that HypG is IND-CCA2 secure,
UF-CMA/SUF-CMA secure, or a secure AEAD in the sense of a published
reduction for its exact composition.  It does not claim NIST Category-5
security, FIPS compliance, certification, or production readiness.  It
does not claim constant-time execution at the machine-code level, and it
does not claim resistance to power, electromagnetic, or fault attacks.
KAT vectors pin behaviour; they are not security evidence.

\subsection{Organisation}

Section~\ref{sec:background} reviews the relevant background.
Section~\ref{sec:principles} states the design constraints.
Section~\ref{sec:principles} describes the architecture and design constraints.
Section~\ref{sec:notation} fixes notation.
Sections~\ref{sec:kem}--\ref{sec:ds} describe the four pillars.
Section~\ref{sec:implementation} discusses implementation assurance.
Section~\ref{sec:evaluation} records the controlled evaluation.
Section~\ref{sec:security} analyses assumptions and open problems.
Section~\ref{sec:validity} discusses threats to validity.
Section~\ref{sec:related} relates HypG to prior work.
Section~\ref{sec:roadmap} gives the path to a defensible profile.

%======================================================================
\section{Background and related work}
\label{sec:background}

\subsection{Post-quantum cryptography}

The NIST post-quantum standardisation process produced ML-KEM
(FIPS~203)~\cite{fips203}, ML-DSA (FIPS~204)~\cite{fips204}, and
SLH-DSA (FIPS~205)~\cite{fips205}.  ML-KEM is a Module-LWE KEM;
ML-DSA is a Module-LWE/Module-SIS signature scheme; SLH-DSA is a
stateless hash-based signature scheme derived from
SPHINCS+~\cite{sphincs}.  HypG does not implement any of these
standards verbatim.  It uses the same broad mathematical families and
some of the same dimension choices, but with different primitives,
encodings, domain separation, and wire formats.

\subsection{Module-LWE}

Module-LWE generalises ring-LWE by working over $R_q^k$.  For a public
matrix $A$ and secret/error vectors $s,e$ drawn from a small
distribution, the public key is $b=A s+e$.  Encryption uses a fresh
small vector $r$ and errors $e_1,e_2$:
\[
  c_1 = A^{T}r+e_1, \qquad
  c_2 = \langle b,r\rangle + e_2 + \Encode(m).
\]
The decryption noise is $\langle e,r\rangle+e_2-\langle s,e_1\rangle$;
correctness requires it to remain inside the decoding region.  The
concrete hardness depends on the ring degree, module rank, modulus,
noise distribution, and number of samples.  The community
lattice-estimator~\cite{estimator} is the standard tool for estimating
the cost of known attacks; it has not been run on HypG's exact
parameters.

\subsection{Hash-based signatures}

SLH-DSA and SPHINCS+ combine a few-time FORS signature with a hypertree
of XMSS/WOTS+ one-time signatures.  The hash function is tweaked by a
public seed and an address so that every hash call is
domain-separated.  HypG follows this architecture with BLAKE3 as the
tweakable hash, a custom 32-byte address, and a deterministic
randomizer.  Determinism is useful for reproducible vectors but removes
fresh-randomness hedging; the security consequences are discussed in
Section~\ref{sec:ds}.

\subsection{Symmetric primitives}

BLAKE3~\cite{blake3} is a tree-structured hash and XOF with a
32-byte keyed mode and a reference \texttt{derive\_key} mode.
ChaCha20~\cite{rfc8439} is a 20-round ARX stream cipher with a 96-bit
nonce and a 32-bit block counter.  HypG uses both as designed: BLAKE3
for hashing, expansion, key derivation, and MACs; ChaCha20 for
confidentiality.  The AEAD is an Encrypt-then-MAC composition rather
than a standardised AEAD algorithm.

%======================================================================
\section{Design principles and architecture}
\label{sec:principles}

\subsection{Five constraints}

Five constraints shape every component.

\begin{enumerate}[leftmargin=*]
\item \textbf{No external dependencies.}  Every dependency edge is a
  workspace path dependency.  The supply-chain surface is therefore the
  Rust toolchain and the repository itself, not a transitive crate
  graph.  \texttt{cargo tree} is part of the release gate.
\item \textbf{No \texttt{unsafe}.}  Every library crate uses
  \texttt{\#![forbid(unsafe\_code)]}.  Arithmetic and conversion paths
  are written to be total where possible.  The cost is scalar code and
  best-effort, rather than volatile, zeroization.
\item \textbf{\texttt{no\_std} cryptographic libraries.}  The core,
  KEM, Symm, AEAD, DS, and hash crates do not require an allocator, a
  filesystem, or an operating system.  Threads, files, and allocation
  live only in \texttt{std} crates.
\item \textbf{Fixed public geometry.}  Protocol loops have public
  bounds.  Secret-dependent branches, secret-dependent array indices,
  and secret-dependent early returns are avoided in the cryptographic
  paths.
\item \textbf{Reproducible evidence.}  Deterministic KATs, a frozen
  checksum inventory, an adversarial harness, a lattice lab,
  cross-architecture equality, and a recorded benchmark are part of the
  artifact.
\end{enumerate}

These are engineering constraints, not security proofs.  In particular,
source-level branchlessness does not imply machine-code constant-time
execution, and safe Rust does not imply resistance to microarchitectural
leakage.

\subsection{Component graph}

The dependency graph is acyclic:
\[
\begin{array}{c}
\texttt{hypg-core}\\[-1mm]
\swarrow\ \downarrow\ \searrow\\[-1mm]
\texttt{hypg-kem}\quad
\texttt{hypg-symm}\to\texttt{hypg-aead}\quad
\texttt{hypg-ds}\to\texttt{hypg-mt}
\end{array}
\]
\texttt{hypg-core} supplies portable BLAKE3, ChaCha20, the CA-PRF
expander, polynomial arithmetic over $\Ring$, and constant-time
helpers.  \texttt{hypg-hash} is a thin API over the core hash
functions.  \texttt{hypg-kat}, \texttt{hypg-bench},
\texttt{hypg-redteam}, \texttt{hypg-lattice}, and \texttt{hypg-cli} are
\texttt{std} harnesses or operator tools.

\subsection{Version identity}

\begin{table}[t]
\centering
\small
\begin{tabular}{@{}llll@{}}
\toprule
Component & Code version & Wire/spec revision & Frozen vectors \\
\midrule
Workspace & 2.0.0 & --- & --- \\
\texttt{hypg-kem} & 2.0.0 & KEM v2.0 & \texttt{kats/v2.0.0.txt} \\
\texttt{hypg-symm} & 2.0.0 & Symm v2.0 & \texttt{kats-symm/v2.0.0.txt} \\
\texttt{hypg-hash} & 2.0.0 & Hash v2.0 & \texttt{kats-hash/v2.0.0.txt} \\
\texttt{hypg-ds} & 2.0.0 & \textbf{DS v2.0.1} & \texttt{kats-ds/v2.0.1.txt} \\
\bottomrule
\end{tabular}
\caption{Version identity.  The literal \texttt{HYP-G/v07/...} strings
are frozen domain separators, not a statement that the parameter set is
v0.7.}
\label{tab:versions}
\end{table}

\subsection{Boundary decisions}

Threading is deliberately excluded from the \texttt{no\_std} crates.
\texttt{hypg-mt} uses \texttt{std::thread::scope} and owns all worker
results; deterministic assembly makes its output byte-identical to the
serial signer.  SIMD intrinsics are excluded by the \texttt{unsafe}
policy and portable SIMD is not stable on the MSRV; the implementation
is scalar and relies on the compiler only for ordinary code generation.

%======================================================================
\section{Notation and mathematical preliminaries}
\label{sec:notation}

\subsection{Bytes and little-endian encoding}

A byte is an element of $\{0,\ldots,255\}$.  For an integer $x$ with
$0\le x<2^t$, $\mathrm{LE}_t(x)$ is the $t/8$-byte encoding with the
least-significant byte first.  All multi-byte protocol integers are
little-endian unless stated otherwise.  Byte ranges are half-open:
$x[a:b]$ denotes bytes $a$ through $b-1$.  The symbol $\concat$ denotes
concatenation.  No implicit padding is inserted.

\subsection{Ring and distributions}

Let
\[
 \Ring=\Z_q[X]/(X^{256}+1),\qquad q=4096=2^{12}.
\]
An element $a\in\Ring$ is represented by $(a_0,\ldots,a_{255})$ with
$a_i\in[0,q)$.  Addition and subtraction are coefficient-wise modulo
$q$; multiplication is negacyclic:
\[
 (a*b)_k=\sum_{\substack{i+j=k\\0\le i,j<256}}a_i b_j
 -\sum_{\substack{i+j=k+256\\0\le i,j<256}}a_i b_j\pmod q.
\]

The distribution $\CBD(2)$ uses four independent uniform bits per
coefficient:
\[
 v=(b_0+b_1)-(b_2+b_3),\qquad b_i\in\bits.
\]
The support is $\{-2,-1,0,1,2\}$ with probabilities
$1/16,4/16,6/16,4/16,1/16$; the mean is zero and the variance is one.

\subsection{Hash notation}

Let $\Hash(x)$ be unkeyed BLAKE3-256, $\Hash_k(x)$ keyed BLAKE3-256 with
32-byte key $k$, and $\XOF(x,L)$ the first $L$ bytes of the BLAKE3 XOF.
Define
\[
\begin{aligned}
 \mathrm{expand}(x,d,L)&=\XOF(x\concat d,L),\\
 \mathrm{kdf}(d,\mathit{parts},L)&=
 \XOF(d\concat\mathit{parts}_0\concat\mathit{parts}_1\concat\cdots,L),\\
 \mathrm{derive}(c,k)&=\mathrm{BLAKE3.derive\_key}(c,k).
\end{aligned}
\]
Neither \texttt{expand} nor \texttt{kdf} inserts length separators; all
protocol callers use fixed-length parts or a single variable-length
part with a distinct domain.

%======================================================================
\section{The CA-PRF expander}
\label{sec:caprf}

\subsection{State and round}

\begin{definition}[CA state and rule]
A CA state is $S\in(\Z_{2^{64}})^{32}$, serialized as 32 little-endian
64-bit words (256 bytes).  A rule is $R\in(\Z_{2^{64}})^8$, serialized
as eight little-endian words (64 bytes).  The public constants are
\[
\begin{aligned}
 \mathrm{RC\_BASE}&=0x9E3779B97F4A7C15,\\
 \mathrm{RC\_STRIDE}&=0xBF58476D1CE4E5B9,\\
 \mathrm{DIFF\_MUL}&=0xD6E8FEB86659FD93.
\end{aligned}
\]
\end{definition}

\begin{definition}[CA round]
For round $t\in[0,32)$ and word $i\in[0,32)$, with all arithmetic
modulo $2^{64}$:
\[
\begin{aligned}
 rc &= \mathrm{RC\_BASE}+(t+1)\cdot\mathrm{RC\_STRIDE},\\
 a &= S_i,\quad b=S_{(i+1)\bmod32},\quad
 c=S_{(i+5)\bmod32},\quad d=S_{(i+17)\bmod32},\\
 rk &= R_{i\bmod8}\oplus rc\oplus(i\cdot\mathrm{RC\_STRIDE}),\\
 t_1 &= a\oplus\mathrm{ROTL}(b,13)\oplus\mathrm{ROTL}(c,7),\\
 t_2 &= (t_1\land d)\oplus((\lnot t_1)\land b),\\
 t_3 &= t_2+(rk\oplus(i+t)),\\
 t_4 &= t_3\oplus\mathrm{ROTL}(t_3,29),\\
 t_5 &= (t_4\cdot\mathrm{DIFF\_MUL})\oplus\mathrm{ROTL}(d,11),\\
 S'_i &= t_5\oplus\mathrm{ROTL}(S_{(i+16)\bmod32},17)\oplus rc.
\end{aligned}
\]
The round is applied 32 times.
\end{definition}

\subsection{CA-PRF definition}

\begin{definition}[CA-PRF]
For a 32-byte seed $x$, public domain $d$, and output length $L$:
\[
\begin{aligned}
 \mathit{init} &= \XOF(x\concat d\concat\texttt{"-init"},256),\\
 \mathit{rule} &= \XOF(x\concat d,64),\\
 S_0 &= \mathrm{LE}^{-1}_{64}(\mathit{init}),\\
 S_{32} &= \mathrm{Round}^{32}(S_0;\mathit{rule}),\\
 \mathit{out} &= \XOF(d\concat\mathrm{LE}_{64}(S_{32})
   \concat d\concat\texttt{"-out"},L).
\end{aligned}
\]
\end{definition}

\begin{remark}[Security role]
The final BLAKE3 XOF is the security-relevant expander.  If BLAKE3 is a
secure XOF or random oracle, the output is pseudorandom regardless of
the CA.  The CA is a deterministic mixing and domain-separation stage,
not an independently proven hardness assumption.
\end{remark}

\begin{lemma}[Prefix consistency]
For fixed $(x,d)$ and $L_1\le L_2$, the first $L_1$ bytes of
$\mathrm{CA\text{-}PRF}(x,d,L_2)$ equal
$\mathrm{CA\text{-}PRF}(x,d,L_1)$.
\end{lemma}

\begin{proof}
The state and all preceding stages are independent of $L$; the final
XOF is prefix-consistent.
\end{proof}

%======================================================================
\section{The Module-LWE KEM}
\label{sec:kem}

\subsection{Parameters and sizes}

\begin{table}[t]
\centering
\small
\begin{tabular}{@{}lrl@{}}
\toprule
Symbol & Value & Meaning \\
\midrule
$n$ & 256 & ring degree \\
$k$ & 4 & module rank; dimension 1024 \\
$q$ & $2^{12}=4096$ & modulus \\
$\eta$ & 2 & $\CBD(2)$ for $s,e,r,e_1,e_2$ \\
$m$ & 32 bytes & FO message; one bit per coefficient \\
$pk$ & 1568 bytes & $\mathit{seed}_A(32)\concat\Pack_{12}(b)(1536)$ \\
$ct$ & 1440 bytes & $\Pack_{10}(c_1)(1280)\concat\Pack_5(c_2)(160)$ \\
$sk$ & 1632 bytes & $\mathit{seed}(32)\concat\mathit{fake}(32)\concat pk$ \\
$ss$ & 32 bytes & XOF output \\
\bottomrule
\end{tabular}
\caption{KEM parameters and sizes.}
\label{tab:kemparams}
\end{table}

The power-of-two modulus differs from ML-KEM-1024's 3329; it enables
branchless mask-only reduction and changes the lattice profile.  ML-KEM
hardness estimates cannot be transferred without a fresh estimator run.

\subsection{Encoding and compression}

\begin{definition}[Encode/Decode]
For a 32-byte message $m$:
\[
 \Encode(m)_i=2048\cdot((m_{\lfloor i/8\rfloor}\gg(i\bmod8))\land1).
\]
For $v\in\Ring$:
\[
 \Decode(v)_i=
 \left(\left(((v_i+1024)\bmod q)\gg11\right)\land1\right),
\]
with output bits packed least-significant first.
\end{definition}

\begin{definition}[Compression]
For $c\in[0,q)$:
\[
 \begin{aligned}
 \Compress_{10}(c)&=((c+2)\gg2)\land0x3FF,\\
 \Decompress_{10}(c)&=(c\land0x3FF)\ll2,\\
 \Compress_5(c)&=((c+64)\gg7)\land0x1F,\\
 \Decompress_5(c)&=(c\land0x1F)\ll7.
 \end{aligned}
\]
\end{definition}

\begin{lemma}[Residual ranges]
For $c\in[0,q)$, let
\[
\begin{aligned}
 \Delta_1&=\Decompress_{10}(\Compress_{10}(c))-c\pmod q,\\
 \Delta_2&=\Decompress_5(\Compress_5(c))-c\pmod q.
\end{aligned}
\]
In the signed range, $\Delta_1\in\{-1,0,1,2\}$ and
$\Delta_2\in\{-63,\ldots,64\}$.
\end{lemma}

\begin{proof}
Write $c=4a+r$ with $0\le r<4$; the 10-bit quotient is $a$ for $r\le1$
and $a+1$ for $r\ge2$, giving $-r$ or $4-r$.  For the 5-bit case,
write $c=128a+r$; the quotient is $a$ for $r<64$ and $a+1$ for
$r\ge64$, giving $-r$ or $128-r$.
\end{proof}

\subsection{Key generation}

\begin{steps}{Algorithm 1: $\KeyGen(\mathit{seed})$}
\item Input: a 32-byte seed with at least 256 bits of entropy.
\item $\mathit{seed}_A\gets\mathrm{CA\text{-}PRF}(\mathit{seed},
  \texttt{"HYP-G/v07/seed-a"},32)$.
\item $A\gets\mathrm{ParseMatrix}(
  \mathrm{CA\text{-}PRF}(\mathit{seed}_A,
  \texttt{"HYP-G/v07/poly-a"},6144))$.
\item $s\gets\mathrm{ParseCBDVector}(
  \mathrm{CA\text{-}PRF}(\mathit{seed},
  \texttt{"HYP-G/v07/sec-s"},512))$.
\item $e\gets\mathrm{ParseCBDVector}(
  \mathrm{CA\text{-}PRF}(\mathit{seed},
  \texttt{"HYP-G/v07/sec-e"},512))$.
\item For $i=0,\ldots,3$:
  $b_i\gets\sum_{j=0}^{3}A_{i,j}*s_j+e_i$ in $\Ring$.
\item $pk\gets\mathit{seed}_A\concat\Pack_{12}(b_0)\concat\cdots
  \concat\Pack_{12}(b_3)$.
\item $\mathit{fake}\gets\mathrm{CA\text{-}PRF}(\mathit{seed},
  \texttt{"HYP-G/v07/fake"},32)$.
\item $sk\gets\mathit{seed}\concat\mathit{fake}\concat pk$.
\item Return $(pk,sk)$; best-effort erase intermediate secrets.
\end{steps}

\subsection{Encapsulation}

\begin{steps}{Algorithm 2: $\Encaps(pk,m)$}
\item Input: a 1568-byte public key and a 32-byte FO message $m$.
  The caller MUST supply a fresh CSPRNG output for $m$ on every call.
\item Parse $pk$ as $(\mathit{seed}_A,b_0,\ldots,b_3)$.
\item $A\gets\mathrm{ParseMatrix}(
  \mathrm{CA\text{-}PRF}(\mathit{seed}_A,
  \texttt{"HYP-G/v07/poly-a"},6144))$.
\item $\mathit{rseed}\gets\Hash(m\concat pk\concat
  \texttt{"HYP-G/v07/r-seed"})$.
\item $r\gets\mathrm{ParseCBDVector}(
  \mathrm{CA\text{-}PRF}(\mathit{rseed},
  \texttt{"HYP-G/v07/rand-r"},512))$.
\item $e_1\gets\mathrm{ParseCBDVector}(
  \mathrm{CA\text{-}PRF}(\mathit{rseed},
  \texttt{"HYP-G/v07/rand-e1"},512))$.
\item $e_2\gets\CBD(\mathrm{CA\text{-}PRF}(\mathit{rseed},
  \texttt{"HYP-G/v07/rand-e2"},128))$.
\item For $i=0,\ldots,3$:
  $c_{1,i}\gets\sum_{j=0}^{3}A_{j,i}*r_j+e_{1,i}$.
\item $c_2\gets\sum_{i=0}^{3}b_i*r_i+e_2+\Encode(m)$.
\item
\[
\begin{aligned}
 ct\gets{}&\Pack_{10}(\Compress_{10}(c_{1,0}))\concat\cdots\\
 &\concat\Pack_{10}(\Compress_{10}(c_{1,3}))
 \concat\Pack_5(\Compress_5(c_2)).
\end{aligned}
\]
\item $ss\gets\XOF(\texttt{"HYP-G/v07/ss"}\concat m\concat ct
  \concat\Hash(pk),32)$.
\item Return $(ct,ss)$.
\end{steps}

\subsection{Decapsulation}

\begin{steps}{Algorithm 3: $\Decaps(sk,ct)$}
\item Parse $sk$ as $(\mathit{seed},\mathit{fake},pk)$.
\item Parse $pk$ and expand $A$ as in encapsulation.
\item $s\gets\mathrm{ParseCBDVector}(
  \mathrm{CA\text{-}PRF}(\mathit{seed},
  \texttt{"HYP-G/v07/sec-s"},512))$.
\item Unpack and decompress $ct$ into $(c_1,c_2)$.
\item $v\gets c_2-\sum_{i=0}^{3}s_i*c_{1,i}$.
\item $m'\gets\Decode(v)$.
\item Recompute $ct'$ exactly as in Algorithm~2 from $(pk,m')$.
\item $\mathit{ok}\gets\mathrm{ct\_eq}(ct',ct)$.
\item $ss_{\mathrm{ok}}\gets\XOF(
  \texttt{"HYP-G/v07/ss"}\concat m'\concat ct
  \concat\Hash(pk),32)$.
\item $ss_{\mathrm{rej}}\gets\XOF(
  \texttt{"HYP-G/v07/reject"}\concat\mathit{fake}\concat ct,32)$.
\item $ss\gets\mathrm{ct\_select}(\mathit{ok},
  ss_{\mathrm{ok}},ss_{\mathrm{rej}})$.
\item Return $ss$; the function never returns an error.
\end{steps}

\subsection{Wire layout}

\begin{table}[t]
\centering
\small
\begin{tabular}{@{}rrl@{}}
\toprule
Offset & Length & Field in $pk$ \\
\midrule
0 & 32 & $\mathit{seed}_A$ \\
32 & 384 & packed $b_0$ \\
416 & 384 & packed $b_1$ \\
800 & 384 & packed $b_2$ \\
1184 & 384 & packed $b_3$ \\
\bottomrule
\end{tabular}
\hfill
\begin{tabular}{@{}rrl@{}}
\toprule
Offset & Length & Field in $ct$ \\
\midrule
0 & 320 & packed compressed $c_{1,0}$ \\
320 & 320 & packed compressed $c_{1,1}$ \\
640 & 320 & packed compressed $c_{1,2}$ \\
960 & 320 & packed compressed $c_{1,3}$ \\
1280 & 160 & packed compressed $c_2$ \\
\bottomrule
\end{tabular}
\caption{KEM public-key and ciphertext layouts.  The secret key is
$\mathit{seed}(32)\concat\mathit{fake}(32)\concat pk(1568)$.}
\label{tab:kemwire}
\end{table}

\subsection{Correctness}

\begin{lemma}[Noise identity]
For a ciphertext produced by Algorithm~2 and unpacked by Algorithm~3,
before compression effects,
\[
 c_2-\langle s,c_1\rangle
 =\Encode(m)+\langle e,r\rangle+e_2-\langle s,e_1\rangle.
\]
\end{lemma}

\begin{proof}
Expand $b=A s+e$ and $c_1=A^T r+e_1$; the mixed terms
$\langle A s,r\rangle$ and $\langle s,A^T r\rangle$ are the same sum
over the same index pairs.
\end{proof}

\begin{lemma}[Compression residual]
Let $\Delta_{1,i}$ be the residual on $c_{1,i}$ and $\Delta_2$ the
residual on $c_2$.  Then
\[
 v'=\Encode(m)+\langle e,r\rangle+e_2-\langle s,e_1\rangle
 +\Delta_2-\langle s,\Delta_1\rangle.
\]
\end{lemma}

\begin{proof}
Substitute $c_{1,i}'=c_{1,i}+\Delta_{1,i}$ and
$c_2'=c_2+\Delta_2$ into $v'=c_2'-\langle s,c_1'\rangle$ and apply
the noise identity.
\end{proof}

\begin{lemma}[Deterministic re-encryption]
If $\Decode(v')=m$, then Algorithm~3 recomputes $ct'=ct$ and selects
$ss_{\mathrm{ok}}$.
\end{lemma}

\begin{proof}
With $m'=m$, every expansion, PKE encryption, compression, and packing
step is deterministic; the fixed-length packed ciphertext equals the
encapsulation value, so the equality test returns $1$.
\end{proof}

\paragraph{Failure probability.}
The working variance model is
\[
 \mathrm{Var}[\langle e,r\rangle]=1024,\quad
 \mathrm{Var}[\langle s,e_1\rangle]=1024,\quad
 \mathrm{Var}[e_2]=1,
\]
\[
 \mathrm{Var}[\langle s,\Delta_1\rangle]=1280,\quad
 \mathrm{Var}[\Delta_2]=1365.25,
\]
for a total of $4694.25$ and $\sigma\approx68.5$.  Rounding up to
$\sigma\approx70$ gives a working margin of approximately
$14.6\sigma$ against the 1024 decision boundary.  This is a heuristic,
not a formal tail bound.  No analytic $2^{-128}$ failure bound is
claimed; the in-tree 32-trial test is a regression tripwire.

\subsection{Intended security argument and gaps}

The intended argument is FO-style: Module-LWE hides $m$ for a passive
adversary; the re-encryption check prevents ciphertext modification; the
random-oracle derivation of $(r,e_1,e_2)$ from $(m,pk)$ prevents
inconsistent randomness; implicit rejection removes a public validity
oracle.  The exact construction is not a verbatim copy of a standard FO
transform~\cite{fo99,hhk17}.  In particular, the accepted secret is
$\KDF(m\concat ct\concat\Hash(pk))$ and the rejection value is
$\KDF(\mathit{fake}\concat ct)$.  No cited theorem covers those exact
choices, and no reduction is proved here.

The lattice parameters have not been run through the community
estimator.  The power-of-two modulus differs from ML-KEM-1024's
3329~\cite{fips203}, so ML-KEM cost estimates cannot be copied.
Consequently the KEM is \emph{not} claimed IND-CCA2 secure; the target
is conditional on the missing analysis.

%======================================================================
\section{Hash, XOF, and KDF layer}
\label{sec:hash}

\subsection{Functions}

The hash layer exposes unkeyed BLAKE3-256, a 64-byte XOF prefix,
keyed BLAKE3, the reference \texttt{derive\_key} construction, an
expander $\XOF(\mathit{seed}\concat\mathit{domain})$, a KDF
$\XOF(\mathit{domain}\concat\mathit{parts}\ldots)$, and a streaming
hasher.  BLAKE3 is used as specified~\cite{blake3}.

\begin{lemma}[XOF prefix property]
$\Hash(x)$ equals the first 32 bytes of $\XOF(x,64)$.
\end{lemma}

\begin{proof}
This is the BLAKE3 root-bytes property, pinned by the repository tests
for empty, short, chunk-edge, and multi-chunk inputs.
\end{proof}

\subsection{Framing requirements}

Neither \texttt{expand} nor \texttt{kdf} inserts separators or lengths.
A caller MUST ensure that its part sequence is injective.  This report
uses fixed-length parts for every KDF call.  A new caller with
variable-length parts MUST add explicit lengths or use
\texttt{derive\_key}.

\subsection{Streaming hasher}

\texttt{Hasher} supports \texttt{new}, \texttt{new\_keyed},
\texttt{update}, \texttt{finalize32}, and \texttt{finalize\_xof}.  The
first 32 bytes of \texttt{finalize\_xof} equal \texttt{finalize32}.
The implementation uses a fixed 54-level chaining-value stack and does
not allocate.

%======================================================================
\section{Authenticated encryption}
\label{sec:aead}

\subsection{One-shot construction}

Let $K$ be a 32-byte master key and $N$ a 12-byte nonce.  Derive
\[
\begin{aligned}
 k_e&=\mathrm{DK}(\texttt{"HYP-G/v07/symm-enc"},K),\\
 k_m&=\mathrm{DK}(\texttt{"HYP-G/v07/symm-mac"},K).
\end{aligned}
\]
Encrypt $C=\mathrm{ChaCha20}(k_e,N,0,M)$ and authenticate
\[
\begin{aligned}
 T=\mathrm{KeyedBLAKE3}\bigl(k_m,\;&
 N\concat AD\concat C\\
 &\concat\mathrm{LE}_{64}(|AD|)
 \concat\mathrm{LE}_{64}(|C|)\bigr).
\end{aligned}
\]
The one-shot opener recomputes the tag and performs a full-scan
constant-time comparison before writing any plaintext.  A tag mismatch
or length mismatch returns \texttt{None} and leaves the plaintext buffer
untouched.

The message limit is exactly $2^{38}$ bytes: $2^{32}$ ChaCha20 blocks of
64 bytes.  Counter exhaustion fails closed rather than wrapping.  The
nonce must be unique per $(K,\text{message})$; the construction has no
nonce-misuse resistance.  ChaCha20 is the IETF 20-round ARX stream
cipher of RFC~8439~\cite{rfc8439}.

\subsection{Streaming front}

The streaming sealer buffers a partial 64-byte keystream block across
calls, feeds ciphertext bytes to a running keyed-BLAKE3 state, and
appends the two little-endian length suffixes at \texttt{finish}.  The
opener feeds ciphertext to the MAC while releasing tentative plaintext;
the plaintext becomes valid only when \texttt{finish} returns
\texttt{Some(())}.  On failure, the assembled plaintext buffer is
zeroized.  Empty chunks, one-byte chunks, block-boundary splits, and
1 MiB streams are tested.

\begin{proposition}[Streaming identity]
For every chunking of the same $(K,N,AD,M)$, the streaming sealer emits
the same $(C,T)$ as one-shot encryption.  For every chunking of the same
$(K,N,AD,C)$, the streaming opener accepts exactly when the one-shot
opener accepts and returns the same plaintext.
\end{proposition}

\begin{proof}[Proof sketch]
The subkeys, initial MAC state, and length counters are identical.  The
keystream state preserves the exact 64-byte block and offset across
calls; by induction on the chunk sequence, the concatenated XOR bytes
equal the one-shot ChaCha20 keystream with counters $0,1,\ldots$.  The
running MAC absorbs exactly $N\concat AD\concat C$ and then the two
length suffixes.  The opener performs the same operations, with the
additional invariant that tentative plaintext is zeroized on failure.
\end{proof}

\subsection{Failure rules}

\begin{table}[t]
\centering
\small
\begin{tabular}{@{}lll@{}}
\toprule
Condition & Sealer & Opener \\
\midrule
length mismatch & erase chunk, poison, $\bot$ & erase chunk, poison, $\bot$ \\
counter exhausted & erase chunk, poison, $\bot$ & erase chunk, poison, $\bot$ \\
length overflow & erase chunk, poison, $\bot$ & erase chunk, poison, $\bot$ \\
finish after poison & erase tag, $\bot$ & erase plaintext, $\bot$ \\
finish length mismatch & --- & erase plaintext, $\bot$ \\
tag mismatch & --- & erase plaintext, $\bot$ \\
success & write tag, $\top$ & validate, $\top$ \\
\bottomrule
\end{tabular}
\caption{Streaming AEAD failure rules.  $\top$ denotes
\texttt{Some(())}; $\bot$ denotes \texttt{None}.}
\label{tab:aeadfail}
\end{table}

A poisoned stream never recovers.  The opener's plaintext is tentative
until \texttt{finish} returns $\top$; callers MUST NOT parse, display,
execute, or transmit it before that point.

\subsection{Security position}

The EtM structure is standard, but the exact ChaCha20 + keyed-BLAKE3
composition is not a standardized AEAD and no reduction for this exact
instantiation is included.  The 32-byte tag is not truncated.  The
construction has no key-commitment guarantee; protocols requiring key
commitment must add it.  Nonce discipline is entirely the caller's
responsibility.

%======================================================================
\section{Hash-based signatures}
\label{sec:ds}

\subsection{Profile and sizes}

The construction follows the SPHINCS+/SLH-DSA family~\cite{sphincs,fips205};
XMSS is standardised in RFC~8391~\cite{rfc8391}.  HypG is not
wire-compatible with either standard.

\begin{table}[t]
\centering
\small
\begin{tabular}{@{}lr@{}}
\toprule
Symbol & Value \\
\midrule
$n$ & 32 bytes (BLAKE3-256) \\
$h$ & 68 \\
$d$ & 17 layers \\
$h'$ & 4 (16 leaves per XMSS tree) \\
$k_{\mathrm{FORS}}$ & 35 \\
$a$ & 9 (512 FORS leaves per tree) \\
$w$ & 16 \\
$\ell$ & 67 (64 message + 3 checksum digits) \\
$pk_{\mathrm{sig}}$ & 64 bytes \\
$sk_{\mathrm{sig}}$ & 128 bytes \\
$\mathit{sig}$ & 49,992 bytes \\
\bottomrule
\end{tabular}
\caption{Signature parameters.}
\label{tab:dsparams}
\end{table}

The sizes are:
\[
\begin{aligned}
 \mathrm{FORS\ unit} &= 32+9\cdot32=320,\\
 \mathrm{FORS\ total} &= 35\cdot320=11{,}200,\\
 \mathrm{XMSS\ layer} &= 8+67\cdot32+4\cdot32=2{,}280,\\
 \mathrm{hypertree} &= 17\cdot2{,}280=38{,}760,\\
 \mathit{sig} &= 32+11{,}200+38{,}760=49{,}992.
\end{aligned}
\]

\subsection{Address structure}

\begin{definition}[ADRS]
An address is 32 bytes:
\[
 \mathrm{layer}(4)\concat\mathrm{tree}(12)\concat
 \mathrm{type}(4)\concat f_{20}(4)\concat f_{24}(4)\concat f_{28}(4),
\]
all integers little-endian.  The 12-byte tree field holds a 96-bit
integer; valid hypertree addresses use at most 68 bits and zero the high
28 bits.
\end{definition}

\begin{table}[t]
\centering
\small
\begin{tabular}{@{}cllll@{}}
\toprule
Type & Name & $f_{20}$ & $f_{24}$ & $f_{28}$ \\
\midrule
0 & \texttt{T\_FORS} & FORS tree & FORS leaf & zero \\
1 & \texttt{T\_WOTS\_HASH} & keypair & chain & absolute step \\
2 & \texttt{T\_WOTS\_PK} & keypair & L-tree height & pair index \\
3 & \texttt{T\_TREE} & FORS tree or zero & node height & node index \\
4 & \texttt{T\_ROOT} & defined, unused & defined, unused & defined, unused \\
\bottomrule
\end{tabular}
\caption{Address types and field assignments.}
\label{tab:adrs}
\end{table}

\subsection{Tweakable hash functions}

All three functions use keyed BLAKE3-256 with the public seed
$\mathit{pubseed}$ as the 32-byte key:
\[
\begin{aligned}
 \PRF(\mathit{pubseed},x,A)
   &=\mathrm{KeyedBLAKE3}(\mathit{pubseed},
     01\concat x\concat A),\\
 F(\mathit{pubseed},A,x)
   &=\mathrm{KeyedBLAKE3}(\mathit{pubseed},
     02\concat A\concat x),\\
 \Hash(\mathit{pubseed},A,L,R)
   &=\mathrm{KeyedBLAKE3}(\mathit{pubseed},
     A\concat L\concat R).
\end{aligned}
\]
The first two inputs are 65 bytes; the third is 96 bytes.  The tags
$01$ and $02$ are the DS 2.0.1 revision.  $\Hash$ is untagged but its
arity and field order are distinct.

\subsection{Key generation}

\begin{steps}{Algorithm 4: $\mathrm{DS.KeyGen}(\mathit{esk})$}
\item $\mathit{pubseed}\gets\mathrm{CA\text{-}PRF}(\mathit{esk},
  \texttt{"HYP-G/v07/ds-pubseed"},32)$.
\item $\mathit{prf}\gets\mathrm{CA\text{-}PRF}(\mathit{esk},
  \texttt{"HYP-G/v07/ds-prf"},32)$.
\item $\mathit{root}\gets\mathrm{HT.Root}(\mathit{pubseed},
  \mathit{esk})$, the top XMSS layer at layer 16, tree 0.
\item $pk\gets\mathit{root}\concat\mathit{pubseed}$.
\item $sk\gets\mathit{esk}\concat\mathit{prf}\concat
  \mathit{pubseed}\concat\mathit{root}$.
\item Return $(pk,sk)$.
\end{steps}

\subsection{Deterministic randomizer and digest}

\begin{steps}{Algorithm 5: $\mathrm{PreSign}(sk,\mathit{msg})$}
\item Parse $(\mathit{esk},\mathit{prf},\mathit{pubseed},
  \mathit{root})$ from $sk$.
\item $\mathit{stream}\gets\mathrm{CA\text{-}PRF}(\mathit{prf},
  \texttt{"HYP-G/v07/ds-addrnd"},64)$.
\item $R\gets\Hash(\mathit{stream}\concat\Hash(\mathit{msg}))$.
\item Compute, for $c\in\{0,1\}$,
  \[
   md_c=\Hash(R\concat\mathit{pubseed}\concat\mathit{root}
   \concat\Hash(\mathit{msg})\concat\mathrm{LE}_{32}(c)
   \concat\texttt{"HYP-G/v07/ds-msghash"}).
  \]
\item Split the first 387 bits of $md=md_0\concat md_1$:
  \[
   \mathit{fors}_t=\sum_{b=0}^{8}2^b\cdot
   \mathrm{bit}(md,9t+b),\quad t=0,\ldots,34,
  \]
  \[
   \mathit{tree}=\sum_{b=0}^{67}2^b\cdot
   \mathrm{bit}(md,315+b),\qquad
   \mathit{leaf}=\sum_{b=0}^{3}2^b\cdot
   \mathrm{bit}(md,383+b).
  \]
\item Return $(R,\mathit{fors},\mathit{tree},\mathit{leaf},
  \mathit{esk},\mathit{pubseed})$.
\end{steps}

There is no fresh per-signature entropy.  Signing is a deterministic
function of $(sk,\mathit{msg})$, so repeated signing returns the same
bytes.  This is intentional for KATs and reproducibility, but it leaks
message equality and removes fresh-randomness fault hedging.  The scheme
is \emph{not} described as freshly hedged.

\subsection{FORS}

\begin{definition}[FORS leaf and node]
For tree $t$ and leaf $l$:
\[
 A=\mathrm{ADRS}(0,0,0,t,l,0),\qquad
 \mathit{leaf}(t,l)=\PRF(\mathit{pubseed},\mathit{esk},A).
\]
For a node at height $h$ and pair index $p$:
\[
 A=\mathrm{ADRS}(0,0,3,t,h,p),\qquad
 \mathit{node}=\Hash(\mathit{pubseed},A,L,R).
\]
\end{definition}

The FORS public key is
\[
 \mathrm{FORS.PK}=\mathrm{KeyedBLAKE3}(\mathit{pubseed},
 A_{\mathrm{root}}\concat\mathit{root}_0\concat\cdots
 \concat\mathit{root}_{34}),
\]
where $A_{\mathrm{root}}=\mathrm{ADRS}(0,0,0,0xFFFFFFFF,0,0)$ and
$\mathit{root}_t$ is the root of tree $t$.  The input is 1152 bytes.
The signature for tree $t$ is the selected leaf plus its nine
authentication nodes.  Verification recomputes each root from the leaf
and path, then compresses the 35 roots.

\begin{lemma}[FORS reconstruction]
Given a valid FORS signature for indices $\mathit{fors}$ and public
seed $\mathit{pubseed}$, $\mathrm{FORS.Verify}$ returns the same
$\mathrm{FORS.PK}$ as FORS public-key generation.
\end{lemma}

\begin{proof}
At height $h$, the verifier uses the current node index
$\lfloor\mathit{fors}_t/2^h\rfloor$, its sibling, and the address
$(t,h,\lfloor\mathit{fors}_t/2^{h+1}\rfloor)$.  The signer stored
exactly that sibling and the same address.  Induction on $h$ therefore
reconstructs each tree root; compression is deterministic.
\end{proof}

\subsection{WOTS+ and L-tree}

\begin{definition}[Digits and checksum]
For a 32-byte message:
\[
 d_i=(msg_{\lfloor i/2\rfloor}\gg4(i\bmod2))\land0x0F,
 \quad i=0,\ldots,63.
\]
The checksum is
\[
 c=\sum_{i=0}^{63}(15-d_i),
\]
encoded as three base-16 digits most-significant first:
$d_{64}=(c\gg8)\land0x0F$,
$d_{65}=(c\gg4)\land0x0F$,
$d_{66}=c\land0x0F$.
\end{definition}

A WOTS+ chain applies 15 tweaked hash steps.  Each step uses a
type-1 address with the keypair index in $f_{20}$, the chain index in
$f_{24}$, and the absolute step in $f_{28}$.  Signing from digit $d$
applies $d$ steps; verification applies $15-d$.  The implementation
always executes 15 calls and selects branchlessly, so the secret digit
does not control the loop trip count.

\begin{lemma}[WOTS chain compatibility]
For $0\le d\le15$, applying $d$ steps and then $15-d$ steps at the
correct absolute addresses equals applying all 15 steps.
\end{lemma}

\begin{proof}
Both compositions apply the same 15 functions with addresses
$0,\ldots,14$; function composition is associative and the segment
boundary at $d$ does not change the addresses.
\end{proof}

The L-tree compresses 67 chain ends.  At height $h$ and pair index $p$:
\[
 A=\mathrm{ADRS}(\mathit{layer},\mathit{tree},2,
 \mathit{keypair},h,p),\qquad
 \mathit{next}_p=\Hash(\mathit{pubseed},A,L,R).
\]
If the width is odd, the final unpaired node is promoted unchanged.  The
widths are $67,34,17,9,5,3,2,1$.

\subsection{XMSS and hypertree}

An XMSS layer has 16 leaves.  For keypair $i$:
\[
 \begin{aligned}
 \mathit{base}&=\mathrm{ADRS}(\mathit{layer},\mathit{tree},0,i,0,0),\\
 \mathit{sk}&=\mathrm{WOTS.GenSK}(\mathit{pubseed},\mathit{esk},
   \mathit{base}),\\
 \mathit{pk}&=\mathrm{WOTS.PKFromSK}(\mathit{pubseed},\mathit{sk},
   \mathit{base}),\\
 \mathit{leaf}&=\mathrm{WOTS.LTree}(\mathit{pubseed},\mathit{pk},
   \mathit{base}).
 \end{aligned}
\]
The Merkle root at height $h$ and node index $p$ uses
\[
 A=\mathrm{ADRS}(\mathit{layer},\mathit{tree},3,0,h,p),
 \qquad
 \mathit{node}=\Hash(\mathit{pubseed},A,L,R).
\]
The XMSS signature is a WOTS+ signature plus four authentication
nodes.  Verification recomputes the leaf and climbs with the same
parity rule.

For layer $i=0,\ldots,16$:
\[
\begin{aligned}
 \mathit{tree}_i&=\mathit{tree\_idx}\gg4(i+1),\\
 \mathit{leaf}_0&=\mathit{leaf},\\
 \mathit{leaf}_i&=(\mathit{tree\_idx}\gg4i)\land15
 \quad(i\ge1).
\end{aligned}
\]
Layer 0 signs the FORS public key; each higher layer signs the root
returned by the layer below.  \texttt{HT.Root} is the XMSS root of layer
16, tree 0.

\begin{lemma}[Hypertree consistency]
For an honestly generated key and signature, hypertree verification
returns \texttt{HT.Root}.
\end{lemma}

\begin{proof}
Induction on the layer index.  The base message is the FORS public key;
at layer $i$, XMSS reconstruction gives the next message; after layer 16
the message is the top root, which equals the public-key root by
construction.
\end{proof}

\subsection{Signing and verification}

\begin{steps}{Algorithm 6: $\mathrm{DS.Sign}(sk,msg)$}
\item $(R,\mathit{fors},\mathit{tree},\mathit{leaf},esk,pubseed)
  \gets\mathrm{PreSign}(sk,msg)$.
\item Compute the 35 FORS units by signing each $\mathit{fors}_t$.
\item Compute the FORS public key from the units.
\item $m_0\gets\mathrm{FORS.PK}$.
\item For $i=0,\ldots,16$:
  \begin{enumerate}[label=\alph*)]
  \item $s_i\gets\mathrm{XMSS.Sign}(\mathit{pubseed},esk,i,
    \mathit{tree}_i,\mathit{leaf}_i,m_i)$.
  \item $m_{i+1}\gets\mathrm{XMSS.Verify}(\mathit{pubseed},i,
    \mathit{tree}_i,\mathit{leaf}_i,m_i,s_i)$.
  \end{enumerate}
\item Serialize $R$, the 35 FORS units, and the 17 layer signatures,
  inserting eight zero bytes before each layer's WOTS+ signature.
\item Return the 49,992-byte signature.
\end{steps}

\begin{steps}{Algorithm 7: $\mathrm{DS.Verify}(pk,msg,sig)$}
\item If $|sig|\ne49{,}992$, return \texttt{false}.
\item Parse $R$, $\mathit{root}$, $\mathit{pubseed}$.
\item Recompute the digest and split into FORS indices, tree index, and
  leaf.
\item Parse the 35 FORS units and compute the FORS public key.
\item For each layer, require the eight reserved bytes to be zero.
\item Climb all 17 XMSS layers in order.
\item Return $\mathrm{ct\_eq}(\mathit{root}',\mathit{root})=1$.
\end{steps}

The signature encoding is
\[
 \mathit{sig}=R(32)\concat
 \mathrm{FORS}_{35}\concat\mathrm{layer}_0\concat\cdots
 \concat\mathrm{layer}_{16}.
\]
Each FORS unit is $32+9\cdot32=320$ bytes.  Each layer is
\[
 \mathrm{zero}(8)\concat\mathrm{WOTS}(67\cdot32)
 \concat\mathrm{auth}(4\cdot32)=2280
\]
bytes.  The total is $32+35\cdot320+17\cdot2280=49{,}992$.
Verification parses every field and rejects nonzero reserved bytes; no
field is silently ignored.

\subsection{DS 2.0.1 revision}

Relative to DS 2.0.0, revision 2.0.1 changes exactly two wire-visible
items:
\begin{enumerate}[leftmargin=*]
\item $\PRF$ prefixes its input with $01$; $F$ prefixes its input with
  $02$.
\item The WOTS+ L-tree promotes an unpaired final node unchanged
  instead of hashing it with itself.
\end{enumerate}
The first removes the 64-byte structural collision between $\PRF$ and
$F$ inputs.  The second follows the FIPS 205 odd-node rule and changes
every WOTS+ L-tree root, every XMSS root, and every signature.  The
frozen file is \texttt{kats-ds/v2.0.1.txt}; the v2.0.0 DS file remains
historical.

\subsection{Security position}

The composition is SLH-DSA-shaped but not FIPS 205 conformant.  It has
no exact UF-CMA/SUF-CMA proof for its BLAKE3 instantiation, address
layout, deterministic randomizer, or reserved-byte encoding.  FORS is a
few-time signature: security degrades as more signatures are produced
under one key.  The implementation enforces no signing limit and keeps
no state.  A deployment profile would need a concrete bound on
signatures per key and a matching multi-target analysis.  The current
report claims only that the implementation computes the specified bytes
and that the deterministic vectors and tamper tests pass.

%======================================================================
\section{Domain registry}
\label{sec:domains}

\begin{longtable}{@{}lll@{}}
\caption{Domain registry.  The \texttt{-init}/\texttt{-out} suffixes are
CA-PRF mechanism, not separate rows.}\\
\toprule
Domain & Construction & Purpose \\
\midrule
\endhead
\texttt{HYP-G/v07/seed-a} & CA-PRF(seed, d, 32) & KEM seed$_A$ \\
\texttt{HYP-G/v07/sec-s} & CA-PRF(seed, d, 512) & KEM secret $s$ \\
\texttt{HYP-G/v07/sec-e} & CA-PRF(seed, d, 512) & KEM error $e$ \\
\texttt{HYP-G/v07/poly-a} & CA-PRF(seed$_A$, d, 6144) & matrix $A$ \\
\texttt{HYP-G/v07/rand-r} & CA-PRF(rseed, d, 512) & FO $r$ \\
\texttt{HYP-G/v07/rand-e1} & CA-PRF(rseed, d, 512) & FO $e_1$ \\
\texttt{HYP-G/v07/rand-e2} & CA-PRF(rseed, d, 128) & FO $e_2$ \\
\texttt{HYP-G/v07/fake} & CA-PRF(seed, d, 32) & rejection key \\
\texttt{HYP-G/v07/r-seed} & BLAKE3(m$\concat$pk$\concat$d) & FO seed \\
\texttt{HYP-G/v07/ss} & XOF(d$\concat$m$\concat$ct$\concat$H(pk),32) & accepted secret \\
\texttt{HYP-G/v07/reject} & XOF(d$\concat$fake$\concat$ct,32) & rejected secret \\
\texttt{HYP-G/v07/symm-enc} & BLAKE3 derive\_key & AEAD enc subkey \\
\texttt{HYP-G/v07/symm-mac} & BLAKE3 derive\_key & AEAD MAC subkey \\
\texttt{HYP-G/v07/ds-addrnd} & CA-PRF(prf, d, 64) & randomizer stream \\
\texttt{HYP-G/v07/ds-msghash} & digest suffix & digest separation \\
\texttt{HYP-G/v07/ds-pubseed} & CA-PRF(esk, d, 32) & public hash key \\
\texttt{HYP-G/v07/ds-prf} & CA-PRF(esk, d, 32) & randomizer seed \\
\texttt{HYP-G/v07/kat} & hash KAT only & never protocol use \\
\bottomrule
\end{longtable}

The retired draft strings \texttt{ds-fors}, \texttt{ds-wots}, and
\texttt{ds-tree} are not used.

%======================================================================
\section{Implementation assurance}
\label{sec:implementation}

\subsection{Memory and allocation discipline}

The cryptographic libraries use fixed-size arrays and caller-provided
buffers.  They do not allocate, read the clock, access the filesystem,
or spawn threads.  The \texttt{std} crates own all OS interaction.  CI
greps the \texttt{no\_std} libraries for \texttt{Vec}, \texttt{Box},
\texttt{alloc::}, \texttt{String}, and \texttt{std::} and fails if any
appears.

\subsection{Secret lifetimes}

Secret-holding types redact \texttt{Debug}, avoid \texttt{Copy} where
duplication would be dangerous, compare in constant time where equality
is exposed, and best-effort zeroize on \texttt{Drop}.  The types include
the KEM secret key and shared secret, the DS secret key and pre-sign
context, the AEAD key, the CA state and rule, the hash state, and
polynomial buffers.  This is best-effort erasure in safe Rust: normal
writes are followed by compiler fences and a black-box use of the
buffer pointer, but safe Rust cannot guarantee volatile erasure.  The
accessors \texttt{seed\_copy}, \texttt{fake\_copy}, and
\texttt{public\_copy} return caller-owned copies that the caller must
erase.

\subsection{Constant-time posture}

The implementation avoids secret-dependent branches and indices in the
cryptographic paths: KEM decapsulation always performs re-encryption and
both KDFs; WOTS+ chains always execute all 15 steps; tag and root
comparisons scan all bytes; implicit-rejection selection is a byte mask.
This is a source-level discipline, not a machine-code guarantee.
Compiler transformations and microarchitectural effects are not
analysed, and no \texttt{dudect}~\cite{dudect} or power/EM/fault
campaign has been performed.

\subsection{Failure semantics}

The KEM never returns an error; invalid ciphertexts yield a
pseudorandom secret.  The one-shot AEAD returns \texttt{None} and leaves
the output untouched on authentication failure.  The streaming opener
marks plaintext as tentative until \texttt{finish}.  Signature
verification returns \texttt{false} on any malformation.  The CLI uses
exit code 2 for usage/input errors and 3 for authentication failures.


%======================================================================
\section{Threat model and failure semantics}
\label{sec:threat}

\subsection{Assets}

\begin{table}[t]
\centering
\small
\begin{tabular}{@{}ll@{}}
\toprule
Asset & Description \\
\midrule
KEM shared secret & 32-byte \texttt{ss} output \\
KEM long-term secret & \texttt{sk = seed} $\concat$ \texttt{fake} $\concat$ \texttt{pk}, especially seed/fake \\
KEM message & 32-byte FO message, exposed as \texttt{coins} \\
AEAD key & 32-byte master key and derived subkeys \\
AEAD plaintext & message and tentative streaming plaintext \\
AEAD associated data & authenticated but public metadata \\
DS master secret & \texttt{esk}, plus derived \texttt{prf} \\
DS public seed/root & public values; integrity matters for verification \\
Signature & 49,992-byte encoding and every parsed field \\
Randomness & key seeds, encapsulation messages, nonces \\
Key copies & stack temporaries, buffers, caller-held copies \\
\bottomrule
\end{tabular}
\caption{Assets considered by the threat model.}
\label{tab:assets}
\end{table}

\subsection{Adversary model}

The intended adversary is computationally bounded and may be classical
or quantum.  The model includes a network adversary that can read,
modify, replay, reorder, truncate, and inject ciphertexts, signatures,
nonces, and associated data; a chosen-ciphertext adversary with a
decapsulation oracle; a chosen-message adversary with a signing oracle;
a malformed-input adversary that supplies arbitrary fixed-length byte
strings; a quantum adversary with Grover search and quantum SVP; a local
software adversary that may attempt memory disclosure or compiler
influence; and a fault adversary that may induce bit flips or transient
faults.  Fault attacks are outside the security claim, although some
fail-closed paths are tested.

The adversary is not assumed to possess the long-term secrets, break
BLAKE3 or ChaCha20, solve Module-LWE, run Shor's algorithm, control the
operating system, or have access to a power/EM laboratory in the
reported evidence.

\subsection{Trust boundaries}

\begin{table}[t]
\centering
\small
\begin{tabular}{@{}p{3.1cm}p{1.5cm}p{7.0cm}@{}}
\toprule
Boundary & Trusted? & Consequence \\
\midrule
Caller CSPRNG & Trusted & weak seeds, coins, or nonces break security \\
Caller key storage & Trusted & library cannot protect keys at rest \\
OS / allocator & Trusted & compromised OS can read or alter memory \\
Compiler / LLVM & Trusted & safe Rust is not a constant-time guarantee \\
Hardware & Trusted & cache, branch predictor, power, and EM are not modelled \\
Protocol framing & Caller & no boundaries, replay protection, or key confirmation \\
Nonce uniqueness & Caller & nonce reuse is catastrophic \\
Signing volume & Caller & FORS is few-time; no limit is enforced \\
Secret copies & Caller & copy accessors can outlive the library wipe \\
\bottomrule
\end{tabular}
\caption{Trust boundaries and caller responsibilities.}
\label{tab:trust}
\end{table}

\subsection{Security goals and non-goals}

The KEM's intended goal is IND-CCA2 with implicit rejection; the AEAD's
intended goals are IND-CPA confidentiality and INT-CTXT integrity under
unique nonces; the DS intended goal is UF-CMA for a bounded signing
volume.  All three are \emph{intended}, not established.

The system does not provide: a complete protocol; key storage, rotation,
revocation, or HSM integration; nonce-misuse resistance; key commitment;
forward secrecy by itself; replay or ordering protection; identity
binding; password-based encryption; protection from a compromised
RNG/OS/compiler/hardware; fault detection; or power/EM/timing
resistance.  It does not provide FIPS validation or NIST certification.

\subsection{Failure semantics}

The KEM never returns an error; invalid ciphertexts yield a
pseudorandom secret.  One-shot AEAD opening returns \texttt{None} and
leaves plaintext untouched on authentication failure.  Streaming AEAD
opening releases only tentative plaintext; failure zeroizes the
assembled buffer.  Signature verification returns \texttt{false} on any
malformation.  The operator CLI uses exit code 2 for usage/input errors
and 3 for authentication failures.

%======================================================================
\section{Evaluation}
\label{sec:evaluation}

\subsection{Method and evidence levels}

We use five evidence levels: E0 source inspection; E1 deterministic
tests and KATs; E2 adversarial harnesses; E3 dynamic/portability
analysis such as Miri and cross-architecture runs; E4 independent
review.  The current artifact has E0--E3 evidence and no E4 evidence.
No result below is an external audit.

The controlled benchmark command was:
\begin{verbatim}
HYPG_BENCH_ITERS=8 \
HYPG_BENCH_DS_ITERS=8 \
HYPG_BENCH_TTEST_N=10000 \
cargo run -q --release -p hypg-bench
\end{verbatim}
The machine was an Intel Core i5-14400F with 16 hardware threads and
31 GiB RAM, running Linux x86-64 with \texttt{rustc 1.97.1}.  The
benchmark warms up once and reports integer averages over eight
iterations; large-message rates use three repetitions of 1 MiB.

\subsection{Correctness and reproducibility}

The debug, release, and MSRV 1.70 test suites each report
\textbf{146 passed, 0 failed, 0 ignored}.  The breakdown is 108 unit
tests, 25 CLI negative integration tests, and 13 doc-tests.  The
red-team harness passes all 14 suites.  The frozen KAT inventory is:

\begin{table}[t]
\centering
\small
\begin{tabular}{@{}lrr@{}}
\toprule
File & Bytes & SHA-256 (prefix) \\
\midrule
\texttt{kats/v2.0.0.txt} & 28,514 & \texttt{3b2dde3d...71d4b} \\
\texttt{kats-symm/v2.0.0.txt} & 3,662 & \texttt{91b7942f...4208} \\
\texttt{kats-hash/v2.0.0.txt} & 923 & \texttt{5fa77803...d9f1} \\
\texttt{kats-ds/v2.0.1.txt} & 201,015 & \texttt{e88c1bdf...c813} \\
\midrule
Total & 234,114 & \\
\bottomrule
\end{tabular}
\caption{Frozen vector inventory.  Complete hashes appear below.}
\label{tab:kats}
\end{table}

The complete SHA-256 values are:
\begin{verbatim}
kats/v2.0.0.txt
3b2dde3d1138c11b54a6880f50ca7b4390071ae53aa94601075f39b2af971d4b
kats-symm/v2.0.0.txt
91b7942ffb4a3c5da8b8b589268af7b7d1594796b428efdc7e26c567c5f04208
kats-hash/v2.0.0.txt
5fa778035ce2e8aa9c89345458ba4da67df4bdbe71a456cbb58026cf9341d9f1
kats-ds/v2.0.1.txt
e88c1bdf54e9390736cd03c1a165616d744a42c48d662880b344c7de3546c813
\end{verbatim}

The x86-64 and \texttt{qemu-aarch64} static-musl KAT outputs are
byte-identical across all 234,114 bytes.  This checks endianness,
integer width, and the DS 2.0.1 revision on an emulated ARM64 target.
It is not native ARM64 validation.

\subsection{Adversarial testing}

The deterministic red-team harness covers garbage inputs, ciphertext
malleability, an $8\times8$ key-mismatch matrix, message extremes, FO
binding, sampler statistics over 524,288 CBD coefficients, failure
sweeps, determinism, AEAD forgery, DS forgery, burst faults, tampered
public keys, padding flips, a panic sweep, secret-key bit flips, and
streaming AEAD games.  A fresh release run with 20,000 trials reported
20,000/20,000 clean decapsulations.  A separate full-bit run flipped all
11,520 bit positions of a test ciphertext; every flip changed the
shared secret.  These are regression and sanity checks, not a
$2^{-128}$ failure proof.

\subsection{Miri and portability}

Miri is an interpreter for Rust's mid-level IR.  Latest recorded runs
on the hardened commit report KEM 9/9 (13,213.79 s), AEAD 6/6
(214.46 s), and DS bounds 3/3 (0.42 s).  Full DS signing is not run
under Miri because the hypertree is hours-long under interpretation.
Miri checks executed paths for some classes of undefined behaviour; it
does not prove their absence.  The libraries also build for
\texttt{aarch64}, \texttt{wasm32}, \texttt{s390x} (check only), and
\texttt{thumbv7em}; KAT and red-team binaries pass under a 1 MB stack.

\subsection{Performance}

\begin{table}[t]
\centering
\small
\begin{tabular}{@{}lr@{}}
\toprule
Operation & Controlled result \\
\midrule
KEM keygen / encaps / decaps & 795 / 931 / 1124 $\mu$s \\
BLAKE3 hash / streaming / XOF (1 MiB) & 842 / 837 / 827 MB/s \\
ChaCha20 XOR (1 MiB) & 773 MB/s \\
AEAD seal / open (1 MiB) & 427 / 434 MB/s \\
DS sign / verify & 72.009 ms / 3.725 ms \\
Threaded DS sign (16 workers) & 18.329 ms, 3.928$\times$ \\
\bottomrule
\end{tabular}
\caption{Controlled release measurements on the record machine.
These are local point estimates, not portable guarantees.}
\label{tab:bench}
\end{table}

The valid-versus-tampered medians were 1004 per-mille for KEM
decapsulation, 602 per-mille for AEAD open, and 994 per-mille for DS
verify.  The Welch t-test over 10,000 interleaved KEM samples had
$\max |t|=0.737$.  These are timing-smoke signals, not a side-channel
proof.

\subsection{Interpretation}

The evaluation establishes that the specified bytes are reproducible,
that the tested failure paths behave as designed, that the KATs survive
debug/release and cross-architecture comparison, and that the scalar
implementation has usable research performance.  It does not establish
cryptographic security, constant-time execution, or production
readiness.


\subsection{Claims and evidence matrix}

\begin{table}[t]
\centering
\small
\begin{tabular}{@{}p{4.7cm}p{4.5cm}p{2.4cm}@{}}
\toprule
Claim & Evidence & Status \\
\midrule
KATs match the implementation & frozen files + tests & verified locally \\
x86-64 and aarch64 KAT bytes agree & cross-arch run & verified locally \\
KEM round-trip and rejection on tested inputs & unit/red-team tests & verified locally \\
KEM IND-CCA2 & no proof; no estimator & not claimed \\
KEM failure $\le 2^{-128}$ & 32-trial tripwire + heuristic & not established \\
AEAD tamper rejection on tested inputs & tests/red-team & verified locally \\
AEAD IND-CPA/INT-CTXT & composition argument only & not proven \\
DS sign/verify and tamper rejection & tests/KATs & verified locally \\
DS UF-CMA/SUF-CMA & no proof; no review & not claimed \\
Memory safety in safe Rust & forbid unsafe; Miri subsets & partially verified \\
Constant-time execution & branchless source patterns & not verified \\
Power/EM/fault resistance & none & not evaluated \\
NIST/FIPS compliance & none & not claimed \\
\bottomrule
\end{tabular}
\caption{Claims and evidence.  E4 (independent review) is absent.}
\label{tab:claims}
\end{table}

%======================================================================
\section{Security analysis and open problems}
\label{sec:security}

\subsection{Assumptions}

The intended security argument depends on: (A1) Module-LWE hardness for
the exact parameters; (A2) BLAKE3 as a collision-resistant hash,
PRF/MAC, and XOF; (A3) ChaCha20 as a secure stream cipher with unique
nonces; (A4) the exact FO-style transform having CCA security under
A1--A2; (A5) random-oracle behaviour of the hash/KDF calls; (A6) a
negligible decapsulation-failure probability; (A7) nonce uniqueness;
(A8) a bounded DS signing volume; (A9) caller-protected secret storage;
and (A10) no machine-code constant-time claim.

\subsection{The exact FO variant}

Standard FO analyses~\cite{fo99,hhk17} do not automatically cover the
chosen shared-secret and rejection derivations.  A proof must adapt the
transform, account for the full public-key binding in $G$, and bound
the re-encryption-collision probability.  This is an open item, not a
footnote.

\subsection{The DS signing volume}

FORS is few-time.  The security of the hypertree depends on how many
messages are signed under one key and on the probability of reusing
FORS leaves or WOTS+ keys.  The current implementation stores no state
and enforces no limit.  A deployment profile needs a concrete bound and
a multi-target analysis; neither exists here.  The deterministic
randomizer also removes fresh-randomness fault hedging.

\subsection{The AEAD composition}

The EtM composition is structurally standard, but the exact
ChaCha20 + keyed-BLAKE3 instantiation has no cited reduction.  Nonce
reuse is catastrophic and caller-controlled.  No key-commitment property
is claimed.

\subsection{Side channels and fault injection}

Source-level branchlessness is not a machine-code guarantee.  The
project has no \texttt{dudect} evaluation, no power/EM laboratory, and
no fault-injection campaign.  The deterministic DS randomizer also
removes fresh-randomness fault hedging.  These are material limitations
for any deployment-oriented claim.

\subsection{Operational assumptions}

The caller must supply a CSPRNG, protect keys in memory, ensure nonce
uniqueness, confirm KEM secrets, bound signature volume, and provide
protocol framing and replay protection.  The libraries provide none of
these services.

%======================================================================
\section{Threats to validity}
\label{sec:validity}

\paragraph{Construct validity.}  KATs and reproducibility checks measure
\emph{behaviour}, not hardness.  A deterministic test can pass while the
construction is insecure.

\paragraph{Internal validity.}  The red-team harness is deterministic and
finite.  It does not cover every input, every protocol composition, or
every compiler transformation.  The 20,000-trial failure sweep is not a
tail bound.

\paragraph{External validity.}  Performance is measured on one CPU,
one toolchain, and one operating system.  Cross-architecture KATs are
emulated on ARM64; native ARM64 and s390x runtime validation are absent.
The results do not establish portable timing behaviour.

\paragraph{Conclusion validity.}  The absence of external audit,
lattice-estimator output, formal proof, and side-channel data prevents
any conclusion about deployable security.

%======================================================================
\section{Related work}
\label{sec:related}

ML-KEM~\cite{fips203}, ML-DSA~\cite{fips204}, and
SLH-DSA~\cite{fips205} are the NIST standards closest to the relevant
families.  HypG shares the Module-LWE and hypertree shapes but is not
wire-compatible or standards-conformant.  SPHINCS+~\cite{sphincs} and
XMSS~\cite{rfc8391} are the main references for the signature
composition.  BLAKE3~\cite{blake3}, ChaCha20~\cite{rfc8439}, and the
Fujisaki--Okamoto literature~\cite{fo99,hhk17} are the symmetric and
transform references.  The lattice-estimator~\cite{estimator} is the
required external tool for the KEM parameters.  The
\texttt{dudect}~\cite{dudect} methodology is the minimum bar for a
timing evaluation.  Rust's Miri~\cite{miri} is used for partial
undefined-behaviour checking.

HypG's contribution is not a new hardness assumption.  It is an
integrated, transparent, reproducible research artifact with an
explicit evidence boundary.

%======================================================================
\section{Road to a defensible profile}
\label{sec:roadmap}

The shortest honest path to a deployable profile is:

\begin{enumerate}[leftmargin=*]
\item Run the community lattice-estimator on the exact
  $(n,k,q,\eta)=(256,4,4096,2)$ parameters, publish classical and
  quantum cost tables, and compare against ML-KEM-1024 without claiming
  equivalence.
\item Derive an analytic decapsulation-failure bound and measure it at
  at least $10^6$ trials.
\item Prove or refute the exact FO-style transform, including the
  accepted and rejected secret derivations.
\item Prove a UF-CMA/SUF-CMA bound for the exact DS composition and
  choose an enforced signing-volume limit.
\item Complete a \texttt{dudect}-style timing evaluation and, if the
  threat model requires it, power/EM/fault evaluation.
\item Add coverage-guided fuzzing and a protocol-level misuse analysis.
\item Commission an independent code audit and publish findings.
\item If FIPS compliance is required, redesign around approved
  primitives and pursue ACVP/CAVP and CMVP validation.
\end{enumerate}

Until then, HypG should remain in research, teaching, reproducibility,
and test environments.

%======================================================================
\section{Conclusion}
\label{sec:conclusion}

HypG v2.0 is a complete, zero-external-dependency, safe-Rust
post-quantum research suite with a documented KEM, hash layer, AEAD, and
stateless signature scheme.  Its evaluation is unusually broad for a
prototype: frozen KATs, cross-architecture equality, 146 tests, a
14-suite adversarial harness, Miri subsets, and controlled performance
measurements.  Its security claims are deliberately narrow: the
artifact implements the specification, the tested paths behave as
specified, and the remaining proof, estimator, side-channel, and audit
work is named explicitly.

That combination---substantial engineering evidence and an explicit
refusal to round it up into a security claim---is the contribution of
this report.

\appendix

%======================================================================
\section{Parameter register}
\label{app:params}

\begin{longtable}{@{}p{3.2cm}p{4.7cm}p{4.1cm}@{}}
\caption{Parameter register.}\\
\toprule
Object & Value & Note \\
\midrule
\endhead
KEM ring & $n=256$, $q=4096$ & $R_q=\Z_q[X]/(X^{256}+1)$ \\
KEM rank & $k=4$ & lattice dimension 1024 \\
KEM noise & $\CBD(2)$ & variance 1, support $\{-2,\ldots,2\}$ \\
KEM sizes & $pk=1568$, $ct=1440$, $sk=1632$, $ss=32$ & bytes \\
DS hash & $n=32$ & BLAKE3-256 \\
DS tree & $h=68$, $d=17$, $h'=4$ & 17 XMSS layers \\
DS FORS & $k=35$, $a=9$ & 35 trees, 512 leaves each \\
DS WOTS+ & $w=16$, $\ell=67$ & 64 message + 3 checksum digits \\
DS sizes & $pk=64$, $sk=128$, $sig=49992$ & bytes \\
AEAD & key 32, nonce 12, tag 32 & bytes \\
AEAD limit & $2^{38}$ & bytes per $(key,nonce)$ \\
CA state & 32 words of 64 bits & 256 bytes \\
CA rounds & 32 & fixed \\
\bottomrule
\end{longtable}

%======================================================================
\section{Domain and API summary}
\label{app:api}

The public KEM API is
\texttt{keygen(seed)}, \texttt{encaps(pk,coins)}, and
\texttt{decaps(sk,ct)}.  The hash API is
\texttt{hash32}, \texttt{hash64}, \texttt{keyed32}, \texttt{derive32},
\texttt{expand}, \texttt{kdf}, and \texttt{Hasher}.  The AEAD API is
\texttt{seal}, \texttt{open}, \texttt{StreamSealer}, and
\texttt{StreamOpener}.  The signature API is
\texttt{ds\_keygen}, \texttt{ds\_sign}, \texttt{ds\_verify},
\texttt{pre\_sign}, \texttt{serialize\_sig}, and \texttt{mt\_sign}.

All fallible APIs return \texttt{Option} or \texttt{bool}; they do not
use exceptions.  Secret-holding types redact \texttt{Debug} and
best-effort zeroize on drop.  The KEM \texttt{coins} argument is the FO
message itself: it must be fresh for every encapsulation.

The exact signatures are:

\begin{verbatim}
keygen(seed: &[u8; 32]) -> (PublicKey, SecretKey)
encaps(pk: &PublicKey, coins: &[u8; 32])
    -> (Ciphertext, SharedSecret)
decaps(sk: &SecretKey, ct: &Ciphertext) -> SharedSecret

hash32(input: &[u8], out: &mut [u8; 32])
hash64(input: &[u8], out: &mut [u8; 64])
keyed32(key: &[u8; 32], input: &[u8], out: &mut [u8; 32])
derive32(context: &str, key_material: &[u8; 32],
         out: &mut [u8; 32])
expand(out: &mut [u8], seed: &[u8; 32], domain: &[u8])
kdf(out: &mut [u8], parts: &[&[u8]], domain: &[u8])

seal(key: &Key, nonce: &Nonce, ad: &[u8], pt: &[u8],
     ct: &mut [u8], tag: &mut [u8; 32]) -> Option<()>
open(key: &Key, nonce: &Nonce, ad: &[u8], ct: &[u8],
     tag: &[u8; 32], pt: &mut [u8]) -> Option<()>
StreamSealer::new/update/finish
StreamOpener::new/update/finish

ds_keygen(esk: &[u8; 32]) -> (SigPk, SigSk)
ds_sign(sk: &SigSk, msg: &[u8], sig: &mut [u8; 49992])
ds_verify(pk: &SigPk, msg: &[u8], sig: &[u8]) -> bool
pre_sign(sk: &SigSk, msg: &[u8]) -> PreSign
serialize_sig(r, fors_sig, ht_sig, sig_out)
mt_sign(sk: &SigSk, msg: &[u8], sig: &mut [u8; SIG_LEN])
\end{verbatim}

%======================================================================
\section{Artifact and reproducibility notes}
\label{app:artifact}

The active documentation is in \texttt{docs/}.  The frozen vectors are
under \texttt{crates/hypg-kat/}.  The controlled benchmark command and
machine details are in \texttt{docs/metrics.md}.  The test matrix and
evidence levels are in \texttt{docs/testplan.md}.  The exact security
non-claims are in \texttt{docs/threat-model.md} and
\texttt{SECURITY.md}.

The core reproducibility commands are:
\begin{verbatim}
cargo test --workspace --locked
cargo test --workspace --locked --release
cargo +1.70 test --workspace --locked
cargo run -q -p hypg-redteam
cargo run -q -p hypg-lattice
cargo test -p hypg-kat --lib --locked
HYPG_BENCH_ITERS=8 \
  HYPG_BENCH_DS_ITERS=8 \
  HYPG_BENCH_TTEST_N=10000 \
  cargo run -q --release -p hypg-bench
\end{verbatim}

The code is licensed under GPL-3.0-or-later.  The report sources are
tracked in \texttt{papers/}; the PDF is built with \texttt{pdflatex}.


%======================================================================
\section{Benchmark detail}
\label{app:bench}

\begin{table}[t]
\centering
\small
\begin{tabular}{@{}llr@{}}
\toprule
Operation & Measurement & Value \\
\midrule
KEM decapsulation & valid median & 1,107,761 ns \\
KEM decapsulation & tampered median & 1,113,174 ns \\
KEM decapsulation & tampered/valid & 1004 per-mille \\
AEAD open 1 KiB & valid median & 3,388 ns \\
AEAD open 1 KiB & forged median & 2,042 ns \\
AEAD open 1 KiB & forged/valid & 602 per-mille \\
DS verify & valid median & 3,697,178 ns \\
DS verify & forged median & 3,678,498 ns \\
DS verify & forged/valid & 994 per-mille \\
\bottomrule
\end{tabular}
\caption{Interleaved valid/tampered medians.  These are timing-smoke
signals, not a leakage proof.}
\label{tab:medians}
\end{table}

\begin{table}[t]
\centering
\small
\begin{tabular}{@{}lr@{}}
\toprule
Welch t-test spot & $t$ \\
\midrule
\texttt{c1[0]} & 0.025 \\
\texttt{c1[320]} & -0.545 \\
\texttt{c2[0]} & 0.737 \\
\texttt{c2[80]} & 0.154 \\
\texttt{ct[last]} & 0.391 \\
\bottomrule
\end{tabular}
\caption{Welch t-statistics for KEM decapsulation, $n=10{,}000$.
The informational bar is $|t|>4.5$; no spot crossed it.  This is not a
\texttt{dudect}-grade evaluation.}
\label{tab:ttest}
\end{table}

\begin{table}[t]
\centering
\small
\begin{tabular}{@{}lr@{}}
\toprule
Micro-profile component & Measurement \\
\midrule
BLAKE3 over 128 B & $\approx187$ ns/hash \\
CA-PRF 1536 B & $\approx4$ $\mu$s \\
BLAKE3 KDF 1536 B & $\approx1$ $\mu$s \\
WOTS+ key (gen + chains + L-tree) & $\approx223$ $\mu$s \\
XMSS layer public key & $\approx3.80$ ms \\
FORS sign & $\approx7.77$ ms \\
FORS verify & $\approx71$ $\mu$s \\
\bottomrule
\end{tabular}
\caption{Optional single-run micro-profile
(\texttt{HYPG\_BENCH\_PROFILE=1}, \texttt{HYPG\_BENCH\_ITERS=4}).
These values are coarse and intended only for cost attribution.}
\label{tab:profile}
\end{table}

The 1 MiB rates were 842 MB/s for BLAKE3 hash, 837 MB/s for streaming
BLAKE3, 827 MB/s for XOF, 773 MB/s for raw ChaCha20 XOR, 427 MB/s for
AEAD seal, and 434 MB/s for AEAD open.  The 1 KiB seal/open averages
were 3/3 $\mu$s.  KEM keygen/encaps/decaps averages were
795/931/1124 $\mu$s; DS sign/verify averages were 72.009/3.725 ms; the
threaded signer averaged 18.329 ms with 16 workers, a 3.928$\times$
speedup.

%======================================================================
\section{Comparison with FIPS 203 and FIPS 205}
\label{app:fips}

\begin{table}[t]
\centering
\small
\begin{tabular}{@{}lll@{}}
\toprule
Property & NIST reference & HypG v2.0.1 \\
\midrule
KEM ring & $q=3329$ (ML-KEM-1024) & $q=4096$ \\
KEM rank & 4 & 4 \\
KEM $pk$ / $ct$ / $sk$ & 1568 / 1568 / 3168 B & 1568 / 1440 / 1632 B \\
KEM compression & $du=11$, $dv=5$ & $c_1=10$, $c_2=5$ \\
KEM hash & SHA-3/SHAKE & BLAKE3 \\
Signature shape & SLH-DSA-256f & same tree dimensions \\
Signature size & 49,856 B & 49,992 B \\
Signature hash & SHA-2/SHAKE & BLAKE3 \\
Signature randomizer & hedged/randomized option & deterministic only \\
NIST status & standard & not standard/validated \\
\bottomrule
\end{tabular}
\caption{Dimensional similarity is not interoperability, compliance, or
a security certification.  The 136-byte signature difference is
$17\times8$ reserved bytes.}
\label{tab:fips}
\end{table}

HypG is not ML-KEM and not SLH-DSA.  The KEM uses a power-of-two modulus,
different compression, a different FO-style secret derivation, and
BLAKE3.  The signature uses a custom 32-byte address, explicit
$\PRF$/$F$ tags, deterministic randomizer derivation, and eight reserved
bytes per layer.  Matching dimensions is not a security certification.

%======================================================================
\section{Symbol glossary}
\label{app:glossary}

\begin{longtable}{@{}ll@{}}
\caption{Symbols used throughout the report.}\\
\toprule
Symbol & Meaning \\
\midrule
\endhead
$\Ring$ & $\Z_q[X]/(X^{256}+1)$, $q=4096$ \\
$s,e$ & KEM long-term secret and error vectors \\
$r,e_1,e_2$ & per-encapsulation FO randomness \\
$m$ & 32-byte FO message; the API's \texttt{coins} \\
$A,b$ & public matrix and public vector $b=As+e$ \\
$pk,ct,sk,ss$ & KEM public key, ciphertext, secret key, shared secret \\
$\CBD(2)$ & centered binomial distribution, variance one \\
$\Hash$ & BLAKE3-256; overloaded as tweakable hash in DS \\
$\XOF$ & BLAKE3 extendable-output function \\
$\KDF$ & domain-separated XOR construction over parts \\
$\PRF,F,H$ & DS tweakable hash functions \\
$R$ & 32-byte DS randomizer \\
$md$ & 64-byte DS message digest \\
$\ADRS$ & 32-byte DS address \\
FORS & few-time forest of 35 height-9 trees \\
WOTS+ & Winternitz one-time signature, $w=16$, $\ell=67$ \\
XMSS & Merkle tree over 16 WOTS+ leaves \\
$h,d,h'$ & total height, layers, per-layer height \\
$sig$ & 49,992-byte signature \\
$\sigma$ & standard deviation; in KEM context, working noise scale \\
\bottomrule
\end{longtable}

\begin{thebibliography}{99}

\bibitem{fips203}
National Institute of Standards and Technology.
\newblock \emph{Module-Lattice-Based Key-Encapsulation Mechanism
Standard}. FIPS 203, August 2024.

\bibitem{fips204}
National Institute of Standards and Technology.
\newblock \emph{Module-Lattice-Based Digital Signature Standard}.
FIPS 204, August 2024.

\bibitem{fips205}
National Institute of Standards and Technology.
\newblock \emph{Stateless Hash-Based Digital Signature Standard}.
FIPS 205, August 2024.

\bibitem{blake3}
J.~O'Connor, J.-P. Aumasson, S.~Neves, and Z.~Wilcox-O'Hearn.
\newblock BLAKE3: one function, fast everywhere.
\newblock Technical report, 2020.
\newblock \url{https://github.com/BLAKE3-team/BLAKE3-specs/}.

\bibitem{rfc8439}
Y.~Nir and A.~Langley.
\newblock ChaCha20 and Poly1305 for IETF Protocols.
\newblock RFC 8439, June 2018.

\bibitem{fo99}
E.~Fujisaki and T.~Okamoto.
\newblock Secure integration of asymmetric and symmetric encryption
schemes.
\newblock In \emph{CRYPTO 1999}, LNCS 1666, pages 537--554, 1999.

\bibitem{hhk17}
D.~Hofheinz, K.~H\"ovelmanns, and E.~Kiltz.
\newblock A modular analysis of the Fujisaki--Okamoto transformation.
\newblock In \emph{TCC 2017}, LNCS 10677, pages 341--371, 2017.

\bibitem{sphincs}
D.~J. Bernstein, A.~H\"ulsing, S.~K\"olbl, R.~Niederhagen,
J.~Rijneveld, and P.~Schwabe.
\newblock The SPHINCS+ signature framework.
\newblock In \emph{ACM CCS 2019}, pages 2129--2146, 2019.

\bibitem{rfc8391}
A.~H\"ulsing, D.~Butin, S.-L. Gazdag, J.~Rijneveld, and A.~Mohaisen.
\newblock XMSS: eXtended Merkle Signature Scheme.
\newblock RFC 8391, May 2018.

\bibitem{estimator}
M.~R. Albrecht, R.~Player, and S.~Scott.
\newblock On the concrete hardness of Learning with Errors.
\newblock \emph{Journal of Mathematical Cryptology}, 9(3):169--203,
2015.

\bibitem{dudect}
O.~Reparaz, J.~Balasch, and I.~Verbauwhede.
\newblock Dude, is my code constant time?
\newblock In \emph{DATE 2017}, pages 1697--1702, 2017.

\bibitem{miri}
The Rust Project.
\newblock Miri: an interpreter for Rust's mid-level intermediate
representation.
\newblock \url{https://github.com/rust-lang/miri}.

\end{thebibliography}

\end{document}
